\documentclass[a4paper]{book}

\setlength{\headheight}{14.5pt}

% --- 1. ENCODAGE ET LANGUE (À mettre en premier) ---
\usepackage[T1]{fontenc}
\usepackage[utf8]{inputenc}
\usepackage[french]{babel}

% --- 2. POLICES ET SYMBOLES ---
% fourier : police élégante pour les maths (lmodern supprimé car doublon)
\usepackage{fourier}
\usepackage{amsfonts, amssymb, amsthm}
\usepackage[right]{eurosym}
\usepackage[np]{numprint}
\usepackage[nointegrals]{wasysym}

% --- 3. GÉOMÉTRIE ET MISE EN PAGE ---
\usepackage[left=1cm,right=1cm,top=1.5cm,bottom=1.5cm]{geometry}
\usepackage[dvipsnames,svgnames]{xcolor}
\usepackage[most,skins,breakable]{tcolorbox}
\usepackage{tagging}

% --- 4. GRAPHISMES ET DIAGRAMMES ---
\usepackage{graphicx}
\usepackage{pgfplots}
\pgfplotsset{compat=1.18}
\usetikzlibrary{babel, shapes, fit} % 'babel' évite les bugs avec TikZ

% --- 5. TABLEAUX ---
\usepackage{array, booktabs, multirow, tabularx, longtable}

% --- 6. STRUCTURE ET TITRES ---
% titlesec est plus moderne que sectsty, on garde uniquement celui-ci
\usepackage{titlesec}
\usepackage{enumitem}
\usepackage[normalem]{ulem} % 'normalem' évite que \emph devienne souligné

% --- 7. HYPERLIENS (toujours en dernier ou presque) ---
\usepackage{hyperref}

\setlength{\parindent}{0pt}
\frenchbsetup{StandardItemLabels=true}

% Le package ultime
\usepackage{ProfCollege}


% =================================================================
% COMMANDES GÉNÉRALES
% =================================================================

% Une ligne de points (usage : \points[15])
\newcommand{\points}[1][1]{%
    \foreach \n in {1,...,#1} { \ldots}%
}

% Plusieurs lignes de points (usage : \lignesPoints{3})
\newcommand{\lignesPoints}[1]{%
    \foreach \n in {1,...,#1} {%
        \dotfill \newline
    }%
}


% =================================================================
% ZONES RÉPONSE
% =================================================================

% Zone réponse universelle (usage : \begin{ZoneReponse}{4cm})
\newtcolorbox{ZoneReponse}[1]{
    enhanced,
    colback=white,
    colframe=gray!50,
    arc=2mm,
    boxrule=0.5mm,
    width=\textwidth,
    height=#1,
    valign=top,
    fonttitle=\bfseries,
    colbacktitle=gray!20,
    coltitle=black,
    attach title to upper,
    after skip=10pt plus 2pt minus 2pt
}


% =================================================================
% COMMANDES MATHÉMATIQUES COLLÈGE
% =================================================================

% Ensembles de nombres
\newcommand{\N}{\ensuremath{\mathbf{N}}}
\newcommand{\Z}{\ensuremath{\mathbf{Z}}}
\newcommand{\D}{\ensuremath{\mathbf{D}}}
\newcommand{\Q}{\ensuremath{\mathbf{Q}}}
\newcommand{\R}{\ensuremath{\mathbf{R}}}

% Valeur absolue
\newcommand{\abs}[1]{\left\lvert #1 \right\rvert}

% Arc géométrique (usage : $\arc{AB}$)
% Note : utilise la fonte tipa pour le symbole arc
\makeatletter
\DeclareFontFamily{U}{tipa}{}
\DeclareFontShape{U}{tipa}{m}{n}{<->tipa10}{}
\newcommand{\arc@char}{{\usefont{U}{tipa}{m}{n}\symbol{62}}}
\newcommand{\arc}[1]{\mathpalette\arc@arc{#1}}
\newcommand{\arc@arc}[2]{%
  \sbox0{$\m@th#1#2$}%
  \vbox{\hbox{\resizebox{\wd0}{\height}{\arc@char}}\nointerlineskip\box0}%
}
\makeatother

% Comparaison : boîte vide entre deux nombres pour écrire le symbole
% Usage : \compare{102}{120}  ou  \compare[0.5]{\np{132999999}}{\np{133000000}}
% Argument optionnel : ignoré (compatibilité ascendante), la boîte s'adapte
\newcommand{\compare}[3][0]{%
    #2 \kern0.4em
    \tikz[baseline=-0.6ex]
        \node[rectangle, draw=gray!60, thick, dotted, minimum height=0.6cm,
              minimum width=0.6cm, inner sep=2pt] {};%
    \kern0.4em #3%
}

% Encadrement : deux boîtes vides autour d'un nombre
% Usage : \encadre{\np{1389}}  ou  \encadre[1.8]{\np{652176}}
% Argument optionnel : largeur des boîtes en cm (défaut 1.5)
\newcommand{\encadre}[2][1.5]{%
    \tikz[baseline=-0.6ex]
        \node[rectangle, draw=gray!60, thick, dotted, minimum height=0.6cm,
              minimum width=#1cm, inner sep=2pt] {};%
    \kern0.3em $<$ #2 $<$ \kern0.3em
    \tikz[baseline=-0.6ex]
        \node[rectangle, draw=gray!60, thick, dotted, minimum height=0.6cm,
              minimum width=#1cm, inner sep=2pt] {};%
}

% Diagramme de proportionnalité simple
% Usage : \diagrammesimple{A}{B}{×3}
\newcommand{\diagrammesimple}[3]{%
\begin{tikzpicture}[baseline=(A.base)]
    \node(A) at (0,0) {#1};
    \node(B) at (5,0) {#2};
    \draw[->, dashed, line width=0.8mm, green!50!black]
        (A) to[bend left=20] node[midway, above=2pt] {#3} (B);
\end{tikzpicture}}

% Diagramme avec flèche aller-retour
% Usage : \diagrammesimplereciproque{A}{B}{×3}{÷3}
\newcommand{\diagrammesimplereciproque}[4]{%
\begin{tikzpicture}[baseline=(A.base)]
    \node(A) at (0,0) {#1};
    \node(B) at (5,0) {#2};
    \draw[->, dashed, line width=0.8mm, green!50!black]
        (A) to[bend left=20] node[midway, above=2pt] {#3} (B);
    \draw[->, dashed, line width=0.8mm, blue!50!black]
        (B) to[bend left=20] node[midway, below=2pt] {#4} (A);
\end{tikzpicture}}
\newcommand{\maClasse}{3e}

%% --- PACKAGES PROPRES AUX COURS ---
\usepackage{fancyhdr}
\usepackage{lastpage}
\usepackage[column=X]{cellspace}

%% --- VARIABLES ---
\newcommand{\monTitre}{Titre du cours} % À redéfinir dans chaque fichier : \renewcommand{\monTitre}{Mon titre}

%% --- FORMAT DES TITRES (titlesec) ---

\titleformat{\chapter}[frame]
{\color{red!85!white}\normalsize}
{\filright\sffamily\Large\enspace Chapitre \thechapter\enspace}
{8pt}
{\color{red!85!white}\rule{0pt}{30pt}\sffamily\Huge\bfseries\filcenter}
\titlespacing{\chapter}{0pc}{0pc}{2pc}

\titleformat{\section}
{\color{green!55!black}\normalfont\Large\bfseries}
{\thesection}
{1em}
{}

% 3. Subsections (Cyan)
% J'ai retiré le \hspace{1em} qui poussait le texte
\titleformat{\subsection}[block]
{\color{cyan!65!black}\normalfont\large\bfseries}
{\thesubsection}
{1em}
{}
% J'ai mis le premier paramètre à 0pt pour supprimer l'indentation de bloc
\titlespacing{\subsection}{0pt}{1.5ex}{1ex}

% 4. Subsubsections (Violet)
\titleformat{\subsubsection}[block]
{\color{violet}\normalfont\normalsize\bfseries}
{\thesubsubsection}
{1em}
{}
\titlespacing{\subsubsection}{0pt}{1ex}{1ex}

%% --- NUMÉROTATION ---
\renewcommand{\thesection}{\Roman{section}.}
\renewcommand{\thesubsection}{\arabic{subsection})}
\renewcommand{\thesubsubsection}{\alph{subsubsection})}
\setcounter{secnumdepth}{3}

%% --- EN-TÊTES ET PIEDS DE PAGE ---
\pagestyle{fancy}
\renewcommand\headrulewidth{1pt}
\fancyhead[LO,RE]{\maClasse}
\fancyhead[RO,LE]{\monTitre}
\renewcommand\footrulewidth{1pt}
\fancyfoot[LO,RE]{Année 2025-2026}
\fancyfoot[RO,LE]{\textbf{Page \thepage/\pageref{LastPage}}}

%% --- RÉGLAGES PARAGRAPHES ---
\setlength{\cellspacebottomlimit}{4pt}
\setlength{\cellspacetoplimit}{4pt}
\setlength{\parindent}{2pt}
\setlength{\tabcolsep}{0.5cm}

%% --- COMPTEURS ET BOITES TCOLORBOX ---

%% --- COMPTEURS AVEC RÉINITIALISATION À CHAQUE CHAPITRE ---

% On ajoute [chapter] à la fin pour que le compteur reparte à 0 à chaque \chapter
\newcounter{dfn}[chapter]
\newcounter{thm}[chapter]
\newcounter{cvn}[chapter]
\newcounter{prp}[chapter]
\newcounter{ex}[chapter]
\newcounter{mtd}[chapter]
\newcounter{rq}[chapter]
\newcounter{demo}[chapter]
\newcounter{appli}[chapter]
\newcounter{exo}[chapter]
\newcounter{lem}[chapter]

%% --- BLOC COMPLET DES BOITES TCOLORBOX ---

% --- DÉFINITIONS (Bleu) ---
\newtcolorbox{definition}{breakable,colback=white!6!white,colframe=blue!85!white,fonttitle=\bfseries,title=Définition \thedfn, code={\stepcounter{dfn}}, arc=3mm}
\newtcolorbox{definitionTitre}[1]{breakable,colback=white!6!white,colframe=blue!85!white,fonttitle=\bfseries,title=Définition \thedfn\ - #1, code={\stepcounter{dfn}}, arc=3mm}
\newtcolorbox{definitions}{breakable,colback=white!6!white,colframe=blue!85!white,fonttitle=\bfseries,title=Définitions \thedfn, code={\stepcounter{dfn}}, arc=3mm}

% --- PROPRIÉTÉS & THÉORÈMES (Rouge/Orange) ---
\newtcolorbox{theoreme}{breakable,colback=red!6!white,colframe=red!85!white,fonttitle=\bfseries,title=Théorème \thethm, code={\stepcounter{thm}}, arc=3mm}
\newtcolorbox{theoremeTitre}[1]{breakable,colback=red!6!white,colframe=red!85!white,fonttitle=\bfseries,title=Théorème \thethm\ - #1, code={\stepcounter{thm}}, arc=3mm}
\newtcolorbox{propriete}{breakable,colback=red!6!white,colframe=red!85!white,fonttitle=\bfseries,title=Propriété \theprp, code={\stepcounter{prp}}, arc=3mm}
\newtcolorbox{proprieteTitre}[1]{breakable,colback=red!6!white,colframe=red!85!white,fonttitle=\bfseries,title=Propriété \theprp\ - #1, code={\stepcounter{prp}}, arc=3mm}
\newtcolorbox{proprietes}{breakable,colback=red!6!white,colframe=red!85!white,fonttitle=\bfseries,title=Propriétés \theprp, code={\stepcounter{prp}}, arc=3mm}
\newtcolorbox{lemme}{breakable,colback=red!6!white,colframe=orange!85!white,fonttitle=\bfseries,title=Lemme \theprp, code={\stepcounter{prp}}, arc=3mm}
\newtcolorbox{convention}{breakable,colback=orange!6!white,colframe=orange!85!white,fonttitle=\bfseries,title=Convention \thecvn, code={\stepcounter{cvn}}, arc=3mm}

% --- EXEMPLES (Vert) — argument optionnel pour le titre ---
% Usage : \begin{exemple}             -> "Exemple 1"
%         \begin{exemple}[Mon titre]  -> "Exemple 1 - Mon titre"
\newtcolorbox{exemple}[1][]{
  enhanced,
  breakable,
  nobeforeafter,
  % --- ON FORCE TOUT À ZÉRO ---
  boxrule=0pt,        % Épaisseur du cadre global à 0
  titlerule=0pt,      % Supprime la ligne sous le titre
  sharp corners,      % Important pour les coupures propres
  % ----------------------------
  colback=white,
  colframe=white,     % On met le cadre en blanc (invisible) au cas où
  frame empty,        % On cache le cadre
  %
  borderline west={0.5mm}{0mm}{green!55!black}, % Seule la barre verte reste
  %
  % Gestion de la coupure
  toprule at break=0pt,
  bottomrule at break=0pt,
  %
  title={Exemple \theex\ifx&#1&\else\ - #1\fi},
  code={\stepcounter{ex}},
  attach boxed title to top left = {xshift = -3mm},
  boxed title style={size=small,colback=green!55!black,colframe=green!55!black, arc=3mm}
}

% --- REMARQUES (Orange) ---
\newtcolorbox{remarque}{enhanced, breakable, frame empty, nobeforeafter, borderline west={0.5mm}{0mm}{orange!85!white}, title=Remarque \therq, colback=white, code={\stepcounter{rq}}, attach boxed title to top left = {xshift = -3mm}, boxed title style={size=small,colback=orange!85!white,colframe=orange!85!white, arc=3mm}}
\newtcolorbox{remarques}{enhanced, breakable, frame empty, nobeforeafter, borderline west={0.5mm}{0mm}{orange!85!white}, title=Remarques \therq, colback=white, code={\stepcounter{rq}}, attach boxed title to top left = {xshift = -3mm}, boxed title style={size=small,colback=orange!85!white,colframe=orange!85!white, arc=3mm}}
\newtcolorbox{remarqueTitre}[1]{enhanced, breakable, frame empty, nobeforeafter, borderline west={0.5mm}{0mm}{orange!85!white}, title=Remarque \therq\ - #1, colback=white, code={\stepcounter{rq}}, attach boxed title to top left = {xshift = -3mm}, boxed title style={size=small,colback=orange!85!white,colframe=orange!85!white, arc=3mm}}

% --- MÉTHODES & APPLICATIONS (Bleu foncé) ---
\newtcolorbox{methode}{enhanced, breakable, frame empty, nobeforeafter, borderline west={0.5mm}{0mm}{blue!75!black}, title=Méthode \themtd, colback=white, code={\stepcounter{mtd}}, attach boxed title to top left = {xshift = -3mm}, boxed title style={size=small,colback=blue!75!black,colframe=blue!75!black, arc=3mm}}
\newtcolorbox{methodes}{enhanced, breakable, frame empty, nobeforeafter, borderline west={0.5mm}{0mm}{blue!75!black}, title=Méthodes \themtd, colback=white, code={\stepcounter{mtd}}, attach boxed title to top left = {xshift = -3mm}, boxed title style={size=small,colback=blue!75!black,colframe=blue!75!black, arc=3mm}}
\newtcolorbox{methodeTitre}[1]{enhanced, breakable, frame empty, nobeforeafter, borderline west={0.5mm}{0mm}{blue!75!black}, title=Méthode \themtd\ - #1, colback=white, code={\stepcounter{mtd}}, attach boxed title to top left = {xshift = -3mm}, boxed title style={size=small,colback=blue!75!black,colframe=blue!75!black, arc=3mm}}
\newtcolorbox{application}{enhanced, breakable, frame empty, nobeforeafter, borderline west={0.5mm}{0mm}{orange!75!black}, title=Application \themtd, colback=white, code={\stepcounter{mtd}}, attach boxed title to top left = {xshift = -3mm}, boxed title style={size=small,colback=orange!75!black,colframe=orange!75!black, arc=3mm}}

% --- DIVERS (Exercices, Python, Conséquences) ---
\newtcolorbox{exercice}{enhanced, breakable, frame empty, nobeforeafter, borderline west={0.5mm}{0mm}{green!55!black}, title=Exercice n°\theexo, colback=white, code={\stepcounter{exo}}, attach boxed title to top left = {xshift = -3mm}, boxed title style={size=small,colback=green!55!black,colframe=green!55!black, arc=3mm}}
\newtcolorbox{exerciceTitre}[1]{enhanced, breakable, frame empty, borderline west={0.5mm}{0mm}{green!55!black}, title=Exercice n°\theexo\ - #1, colback=white, code={\stepcounter{exo}}, attach boxed title to top left = {xshift = -3mm}, boxed title style={size=small,colback=green!55!black,colframe=green!55!black, arc=3mm}}
\newtcolorbox{consequences}{enhanced, breakable, frame empty, nobeforeafter, borderline west={0.5mm}{0mm}{ProcessBlue}, title=Conséquences, colback=white, attach boxed title to top left = {xshift = -3mm}, boxed title style={size=small,colback=ProcessBlue,colframe=ProcessBlue, arc=3mm}}
\newtcolorbox{programme}{enhanced, breakable, frame empty, nobeforeafter, borderline west={0.5mm}{0mm}{gray!55!black}, title=Programme Python, colback=white, attach boxed title to top left = {xshift = -3mm}, boxed title style={size=small,colback=gray!55!black,colframe=gray!55!black, arc=3mm}}


% Commande pour les démonstrations (incrémentation automatique)
\newcommand{\demonstration}{%
  \stepcounter{demo}%
  \tcbox[enhanced, 
         on line, 
         arc=2mm, 
         colback=gray!10!white, 
         colframe=gray!50!black, 
         size=small, 
         fontupper=\bfseries]
         {Démonstration \thedemo}%
  \enspace % Ajoute un petit espace après la boîte
}

%% --- FIN DU BLOC DES BOITES ---

% Activité (Numéro, Titre, Chapitre)
\newcommand{\enteteActivite}[3]{
    \noindent
    \begin{tabularx}{\textwidth}{@{}p{2cm} >{\centering\arraybackslash}X p{2cm}@{}}
        \textbf{Activité n°#1} & \textbf{#2} & \textbf{Chapitre #3} \\
    \end{tabularx}
}


% --- Commande pour les blocs mémo personnalisés ---
% Argument 1 : Couleur du cadre
% Argument 2 : Titre du bloc
% Argument 3 : Texte/Contenu
\newcommand{\blocMemo}[3]{%
    \begin{tcolorbox}[
        enhanced,
        breakable,
        title=#2,
        colback=#1!5,          % Fond très clair
        colframe=#1!75!black,  % Bordure
        fonttitle=\bfseries\sffamily,
        colbacktitle=#1!75!black,
        attach boxed title to top left={xshift=3mm, yshift=-3mm},
        boxed title style={size=small, arc=2mm, colframe=#1!75!black},
        arc=3mm,               % Bel arrondi
        boxrule=0.6mm,         % Trait un peu plus épais
        after skip=15pt,
        before skip=10pt
    ]
        \vspace{2mm} % Petit espace pour ne pas coller au titre boxé
        #3
    \end{tcolorbox}
}


% --- Nouvel environnement Culture Maths ---
\newtcolorbox{cultureMaths}{
    enhanced, 
    breakable, 
    frame empty, 
    nobeforeafter, 
    borderline west={0.5mm}{0mm}{gray!60!black}, % Bordure grise comme ton "programme"
    title=Le saviez-vous ?, 
    colback=white, 
    attach boxed title to top left = {xshift = -3mm}, 
    boxed title style={size=small,colback=gray!60!black,colframe=gray!60!black, arc=3mm},
    coltitle=white,
    fonttitle=\bfseries\sffamily,
    after skip=15pt,
    before skip=10pt
}

\begin{document}

\setcounter{chapter}{-1}
\renewcommand{\monTitre}{Introduction aux notions de fonctions}
\chapter{\monTitre}
\thispagestyle{fancy}

    \section{Qu'est-ce qu'une fonction ?}
    
    \begin{definition}
        Une \textbf{fonction} est comme une \textbf{machine} mathématique. 
        On lui donne un nombre au départ (le nombre $x$), elle lui fait subir des calculs, et elle ressort un \textbf{unique} résultat.
        
        Pour pouvoir définir une \textbf{une fonction} il y a donc une contrainte forte : 
        
        \begin{center}
            Un nombre de départ ne peut avoir qu'un unique résultat.
        \end{center}
    \end{definition}
    
    \begin{exercice}[Le test du diagramme]
        Parmi les diagrammes suivants, lesquels peuvent représenter une fonction ? Justifier.
        
        \medskip
        \begin{center}
        \begin{tikzpicture}[scale=0.75]
            % Diagramme A
            \node at (1,2.5) {\textbf{A}};
            \draw (0,0) ellipse (0.6 and 1.5);
            \draw (2,0) ellipse (0.6 and 1.5);
            \node (a1) at (0,0.8) {1}; \node (a2) at (0,-0.8) {2};
            \node (b1) at (2,0.8) {5}; \node (b2) at (2,-0.8) {10};
            \draw[->,>=latex] (a1) -- (b1);
            \draw[->,>=latex] (a1) -- (b2);
            \draw[->,>=latex] (a2) -- (b2);
            
            % Diagramme B
            \begin{scope}[xshift=4.5cm]
            \node at (1,2.5) {\textbf{B}};
            \draw (0,0) ellipse (0.6 and 1.5);
            \draw (2,0) ellipse (0.6 and 1.5);
            \node (a1) at (0,0.8) {1}; \node (a2) at (0,-0.8) {2};
            \node (b1) at (2,0) {7};
            \draw[->,>=latex] (a1) -- (b1);
            \draw[->,>=latex] (a2) -- (b1);
            \end{scope}
    
            % Diagramme C
            \begin{scope}[xshift=9cm]
            \node at (1,2.5) {\textbf{C}};
            \draw (0,0) ellipse (0.6 and 1.5);
            \draw (2,0) ellipse (0.6 and 1.5);
            \node (a1) at (0,1) {4}; \node (a2) at (0,0) {8}; \node (a3) at (0,-1) {9};
            \node (b1) at (2,1) {2}; \node (b2) at (2,-1) {3};
            \draw[->,>=latex] (a1) -- (b1);
            \draw[->,>=latex] (a2) -- (b1);
            \draw[->,>=latex] (a3) -- (b2);
            \end{scope}
    
            % Diagramme D (VIDE)
            \begin{scope}[xshift=13.5cm]
            \node at (1,2.5) {\textbf{D (À toi !)}};
            \draw (0,0) ellipse (0.6 and 1.5);
            \draw (2,0) ellipse (0.6 and 1.5);
            \node at (0,-1.8) {\textit{Départ}};
            \node at (2,-1.8) {\textit{Arrivée}};
            \end{scope}
        \end{tikzpicture}
        \end{center}
    
        \begin{enumerate}
            \item Les fonctions sont les diagrammes : \dotfill
            \item Justification : \begin{ZoneReponse}{1.8cm}\end{ZoneReponse}
            \item \textbf{Défi :} Dans le cadre \textbf{D}, invente ton propre diagramme (nombres et flèches) pour qu'il \textbf{représente bien une fonction}.
        \end{enumerate}
    \end{exercice}

    
\section{La machine à calculer}
    
    \begin{exercice}[Fonctionnement interne]
        On considère une fonction $f$ représentée par la machine suivante :
        \begin{center}
            \fbox{Nombre de départ} $\longrightarrow$ \fbox{$\times 3$} $\longrightarrow$ \fbox{$- 5$} $\longrightarrow$ \fbox{Résultat}
        \end{center}
        
        \begin{enumerate}
            \item Si on choisit $4$ au départ, quel est le résultat ? \dotfill
            \item Si on choisit $-2$ au départ, quel est le résultat ? \dotfill
            \item Quel nombre faut-il choisir au départ pour obtenir $16$ ? \dotfill
            \item Quel nombre faut-il choisir au départ pour obtenir $1$ ? \dotfill
            \item \textbf{Notation :} Si on appelle $x$ le nombre de départ, exprime le résultat $f(x)$ en fonction de $x$. \\ $f(x) = $ \dotfill
        \end{enumerate}
    \end{exercice}
    
\section{Lecture graphique (Le "Montagnes Russes")}
    
    \begin{exercice}[Analyse d'une courbe]
        Voici la représentation graphique d'une fonction $g$.
        
        \begin{center}
        \begin{tikzpicture}[scale=0.8, grid help/.style={help lines, dashed, gray!40}]
            \draw[grid help] (-5,-3) grid (6,4);
            \draw[thick,->] (-5,0) -- (6,0) node[right] {$x$};
            \draw[thick,->] (0,-3) -- (0,4) node[above] {$y$};
            \foreach \x in {-4,-3,...,5} \draw (\x,0.1) -- (\x,-0.1) node[below] {\footnotesize \x};
            \foreach \y in {-2,-1,1,2,3} \draw (0.1,\y) -- (-0.1,\y) node[left] {\footnotesize \y};
            
            \draw[very thick, blue] plot [smooth, tension=0.8] coordinates {(-4,3) (-2,1) (-1,2) (0,0) (1,-2) (3,-1) (5,3)};
            \node[blue] at (4.5,1.5) {$\mathcal{C}_g$};
        \end{tikzpicture}
        \end{center}
    
    
        En utilisant le graphique, répondre aux questions suivantes :
        \begin{enumerate}
            \item Quel est le nombre associé au nombre de départ $x = -2$ par cette fonction ? \dotfill
            \item Pour quels nombres de départ a-t-on un résultat égal à $0$ ? \dotfill
            \item Quel est le résultat maximum atteint par cette machine ? \dotfill
            \item Citer un nombre qui n'a \textbf{aucun résultat} possible par cette fonction. \dotfill
        \end{enumerate}
    \end{exercice}

    
\section{Notations et Vocabulaire}

    \blocMemo{RoyalBlue}{Le langage des fonctions (À retenir absolument)}{
        Soit $f$ une fonction qui, à un nombre $x$, associe un résultat $y$. 
        
        \medskip
        \begin{itemize}
            \item \textbf{La notation :} On écrit $f : x \longmapsto 3x - 5$ ou plus simplement $f(x) = 3x - 5$ (se lit \og $f$ de $x$ égale $3x -5$ \fg ).
            \item \textbf{L'Antécédent :} C'est le nombre de \textbf{départ}; ici $x$. 
            \item \textbf{L'Image :} C'est le nombre d' \textbf{arrivée}, ici le résultat $f(x)$.
        \end{itemize}
        
        \begin{center}
            \begin{tikzpicture}
                \node (x) at (0,0) {\Large $x$};
                \node (f) [draw, rectangle, minimum width=2cm, minimum height=1cm, fill=gray!10] at (3,0) {\Large $f$};
                \node (fx) at (6,0) {\Large $f(x)$};
                \draw[->, thick] (x) -- (f);
                \draw[->, thick] (f) -- (fx);
                \node[below=0.3cm of x, color=RoyalBlue] {\textbf{Antécédent}};
                \node[below=0.3cm of fx, color=OrangeRed] {\textbf{Image}};
            \end{tikzpicture}
        \end{center}
        
        \textit{Exemple : Si $f(2) = 7$, on dit que \textbf{7 est l'image de 2}, et que \textbf{2 est un antécédent de 7}.}
    }
    

    \begin{exercice}[Le jeu des traductions]
        Compléter les correspondances suivantes (vocabulaire et notation) :
    
        \medskip
        \begin{itemize}[label={$\bullet$}, itemsep=0.8cm]
            \item \textbf{Notation :} $f(5) = 12$ \\[1ex]
                  \textbf{Phrase :} \dotfill
            
            \item \textbf{Phrase :} L'image de $-3$ par la fonction $g$ est $10$. \\[1ex]
                  \textbf{Notation :} \dotfill

\newpage
                  
            \item \textbf{Notation :} $h : x \longmapsto 4x$ \\[1ex]
                  \textbf{Phrase :} \dotfill
                  
            \item \textbf{Phrase :} Un antécédent de $8$ par la fonction $k$ est $2$. \\[1ex]
                  \textbf{Notation :} \dotfill
                  
            \item \textbf{Notation :} $f(0) = -4$ \\[1ex]
                  \textbf{Phrase :} \dotfill
        \end{itemize}
    \end{exercice}

    
    \begin{exercice}[Vrai ou Faux ?]
        On considère une fonction $p$ telle que $p(4) = -1$.
        \begin{enumerate}
            \item L'image de $4$ est $-1$. \quad $\square$ Vrai \quad $\square$ Faux
            \item $4$ est un antécédent de $-1$. \quad $\square$ Vrai \quad $\square$ Faux
            \item $-1$ est l'antécédent de $4$. \quad $\square$ Vrai \quad $\square$ Faux
            \item L'image de $-1$ est $4$. \quad $\square$ Vrai \quad $\square$ Faux
        \end{enumerate}
    \end{exercice}


    \begin{exercice}[Le sens inverse]
        On reprend la courbe $\mathcal{C}_g$ de l'exercice de la page précédente.
        \begin{enumerate}
            \item Trouve tous les \textbf{antécédents} de $1$ par la fonction $g$.
            \item Combien le nombre $-2$ possède-t-il d'antécédents ?
            \item Existe-t-il un nombre qui a \textbf{3 antécédents} ? Si oui, lequel ?
        \end{enumerate}
        \begin{ZoneReponse}{2.5cm}\end{ZoneReponse}
    \end{exercice}


    \begin{methode}[Calculer une image par une formule]
        Soit la fonction $f$ définie par l'expression : $f(x) = x^2 - 5x + 2$.
        
        \textbf{Question :} Quelle est l'image de $3$ par la fonction $f$ ?
        
        \textbf{Réponse :} On remplace $x$ par $3$ dans l'expression :
        \[ f(3) = 3^2 - 5 \times 3 + 2 \]
        \[ f(3) = 9 - 15 + 2 = -4 \]
        L'image de $3$ par $f$ est donc $-4$.
    \end{methode}
    
    
    \begin{exercice}
        On considère la fonction $f$ définie par $f(x) = x^2 - 5x + 2$.
        \begin{enumerate}
            \item Calcule l'image de $10$ par $f$ : \\
            $f(10) = $ \dotfill
            
            \item Calcule $f(-2)$ (Attention aux parenthèses !) : \\
            $f(-2) = $ \dotfill
        \end{enumerate}

        \medskip
        \textbf{Défi : Recherche d'antécédent} \\
        On considère maintenant une autre fonction $h$ définie par $h(x) = 3x - 4$.
        \begin{enumerate}[resume]
            \item Quel est l'antécédent de $11$ par la fonction $h$ ? \\
            (Indication : Tu dois chercher la valeur de $x$ pour que $3x - 4 = 11$) \\
        \end{enumerate}
    \end{exercice}


\end{document}